\documentclass[11pt]{article}
\usepackage[letterpaper,margin=0.9in]{geometry}
\usepackage[T1]{fontenc}
\usepackage{lmodern}
\usepackage{amsmath,amssymb,amsthm,mathtools,mathrsfs,bm,microtype}
\usepackage{enumitem,longtable,booktabs,array,tabularx,graphicx}
\usepackage{xparse,url,needspace}
\usepackage[hidelinks,hypertexnames=false]{hyperref}
\numberwithin{equation}{section}
\newtheorem{theorem}{Theorem}[section]
\newtheorem{lemma}[theorem]{Lemma}
\newtheorem{proposition}[theorem]{Proposition}
\newtheorem{corollary}[theorem]{Corollary}
\theoremstyle{definition}
\newtheorem{definition}[theorem]{Definition}
\theoremstyle{remark}
\newtheorem{remark}[theorem]{Remark}
\newtheorem{notation}[theorem]{Notation}
\theoremstyle{plain}
\newtheorem{targettheorem}{Theorem}
\NewDocumentEnvironment{namedtheorem}{m o}
 {\renewcommand{\thetargettheorem}{#1}\IfNoValueTF{#2}
    {\begin{targettheorem}}{\begin{targettheorem}[#2]}}
 {\end{targettheorem}}
\newcommand{\R}{\mathbb R}
\newcommand{\C}{\mathbb C}
\newcommand{\Lsig}{L^2_\sigma(\R^3)}
\newcommand{\eps}{\varepsilon}
\newcommand{\ip}[2]{\langle#1,#2\rangle}
\DeclareMathOperator{\supp}{supp}
\DeclareMathOperator{\ran}{ran}
\DeclareMathOperator{\dist}{dist}
\DeclareMathOperator{\curl}{curl}
\DeclareMathOperator{\esssup}{ess\,sup}
\allowdisplaybreaks[2]
\setcounter{secnumdepth}{3}
\setcounter{tocdepth}{1}
\setlength{\emergencystretch}{3em}
\Urlmuskip=0mu plus 1mu
\raggedbottom
\newcommand{\CompleteCandidate}{1}
\newcommand{\ProofModule}[2]{%
  \IfFileExists{#1}{\input{#1}}{%
    \par\medskip\noindent\textbf{Unproved dependency.} #2
    The proof of this step is not supplied in this draft.\par\medskip
    \PackageWarningNoLine{complete-candidate}{Missing proof module: #1}}}
\title{Non-Uniqueness of Leray--Hopf Solutions to the Unforced\\
Three-Dimensional Navier--Stokes Equations}
\author{Tristan Buckmaster \and Matei P. Coiculescu \and Chen Li}
\date{November 2026}
\begin{document}
\maketitle
\begin{abstract}
We prove that the incompressible Navier--Stokes equations in $\R^3$,
without external force, admit a continuum of distinct suitable weak
solutions in the Leray--Hopf class with the same initial datum. The datum is compactly supported,
axisymmetric without swirl, and smooth away from the origin; it lies in
$L^{3,\infty}$ and in $L^q$ for every $q<3$, but not in $L^3$. The
solutions are axisymmetric without swirl and smooth at positive times,
and they satisfy the local and global energy equalities. Any two
of them separate in $L^2$ at the rate $t^{1/4+\sigma}$, for a fixed
$0<\sigma<1/4$. The proof follows the program of Jia and \v Sver\'ak,
with a discretely self-similar background in place of a self-similar one.
The background is an exact unforced solution built from a train of thin
vortex rings, which we obtain by lifting a periodic row of planar vortices
with Gaussian stream function. We construct a growing Floquet mode of its
linearization directly, from the angular-two instability of the planar
vortex, and conclude by a nonlinear correction and a localization of the
datum.
\end{abstract}

\tableofcontents

\section{Introduction}
\label{sec:intro}

We consider the incompressible Navier--Stokes equations in $\R^3$
without external force:
\begin{equation}\label{eq:ns}
 \partial_tu+\operatorname{div}(u\otimes u)-\Delta u+\nabla p=0,
 \qquad \operatorname{div}u=0,\qquad u(0)=u_0.
\end{equation}
Here $u$ is the velocity, $p$ is the pressure, and the viscosity is
normalized to one. For every divergence-free $u_0\in L^2(\R^3)$,
Leray~\cite{leray1934} constructed a global weak solution that satisfies
the energy inequality
\[
 \frac12\|u(t)\|_2^2+\int_0^t\|\nabla u(s)\|_2^2\,ds
 \le\frac12\|u_0\|_2^2,
\]
and Hopf~\cite{hopf1951} gave the corresponding construction in bounded
domains. We call such solutions Leray--Hopf solutions. Whether they are
uniquely determined by their datum is a basic question that goes back to
Leray's work. In this paper we prove that, without force, uniqueness fails
in the axisymmetric class without swirl: Theorem~A below gives a compactly
supported datum from which a continuum of distinct Leray--Hopf solutions
emanates. The proof is analytic. Recent results on the same question are
discussed in Section~\ref{sec:intro-nonuniqueness}.

\subsection{Uniqueness and regularity}
\label{sec:intro-uniqueness}

Uniqueness holds as long as one of the solutions is regular. If a
Leray--Hopf solution lies in $L^r(0,T;L^q)$ with $2/r+3/q=1$ and
$3<q\le\infty$, then every Leray--Hopf solution with the same datum
coincides with it on $[0,T]$. This weak--strong uniqueness theorem is due
to Prodi~\cite{prodi1959}, Serrin~\cite{serrin1963} and
Ladyzhenskaya~\cite{ladyzhenskaya1967}; Kozono and
Sohr~\cite{kozono-sohr1996} treated the endpoint $L^\infty(0,T;L^3)$.
Such regular solutions exist for data in the critical space $L^3$:
Kato~\cite{kato1984} constructed solutions for data in $L^3$, local in
time for large data and global for small data. Koch and Tataru~\cite{koch-tataru2001} obtained global
solutions for small data in $BMO^{-1}$. Escauriaza, Seregin and
\v Sver\'ak~\cite{ess2003} proved that Leray--Hopf solutions bounded in
$L^\infty(0,T;L^3)$ are smooth. For a general Leray--Hopf solution only
partial regularity is known. Scheffer~\cite{scheffer1976} bounded the
size of the singular set at the first singular time, and Caffarelli, Kohn
and Nirenberg~\cite{ckn1982} proved that suitable weak solutions, those
satisfying the local energy inequality, are smooth outside a closed set of
zero one-dimensional parabolic Hausdorff measure.
We refer to the monograph~\cite{lemarie2016} for these results.

The critical spaces are those invariant under the scaling
$u\mapsto\lambda u(\lambda x,\lambda^2t)$ of~\eqref{eq:ns}. The datum of
Theorem~A lies in the weak Lebesgue space $L^{3,\infty}$ but not in
$L^3$; near the origin it is of size $|x|^{-1}$. For such data the theory
is incomplete. Seregin and \v Sver\'ak~\cite{seregin-sverak2017}
studied global weak solutions with large data in $L^3$, and Barker,
Seregin and \v Sver\'ak~\cite{bss2018} constructed, for data in
$L^{3,\infty}$, global weak solutions that are stable under weak-$*$
convergence of the data; their uniqueness is not known.
For a scale-invariant datum it is natural to look for a self-similar
solution $u(x,t)=t^{-1/2}U(x/\sqrt t)$. Jia and \v Sver\'ak~\cite{jia-sverak2014} proved
that every scale-invariant datum that is smooth away from the origin,
of any size, gives rise to at least one forward self-similar solution.
Discretely self-similar data, invariant only under
$u_0(x)\mapsto\lambda u_0(\lambda x)$ for one $\lambda>1$, also admit
discretely self-similar solutions: Tsai~\cite{tsai2014} constructed them
for $\lambda$ close to one, and Bradshaw and Tsai~\cite{bradshaw-tsai2017}
and Chae and Wolf~\cite{chae-wolf2018} for every $\lambda>1$. Our solutions are built around a
discretely self-similar solution of this kind.

In the axisymmetric class without swirl the theory is more complete.
Ladyzhenskaya~\cite{ladyzhenskaya1968} and Ukhovskii and
Yudovich~\cite{ukhovskii-yudovich1968} proved that smooth, sufficiently
decaying axisymmetric data without swirl give global smooth solutions, and
Leonardi, M\'alek, Ne\v cas and Pokorn\'y~\cite{lmnp1999} relaxed the
regularity of the data. Feng and \v Sver\'ak~\cite{feng-sverak2015}
constructed solutions whose initial vorticity is a vortex ring, a measure
supported on circles, and Gallay and \v Sver\'ak~\cite{gallay-sverak2019}
proved that the solution emanating from a circular vortex filament is
unique. The datum of Theorem~A is axisymmetric without swirl and smooth
away from the origin. Theorem~A therefore shows that a singularity at one
point, of the critical size $|x|^{-1}$, is enough for uniqueness to fail
in this class.

\subsection{Non-uniqueness}
\label{sec:intro-nonuniqueness}

Two methods have produced non-unique solutions of the equations of fluid
mechanics. The first is convex integration.
Buckmaster and Vicol~\cite{buckmaster-vicol2019} proved that weak solutions of the
Navier--Stokes equations with finite kinetic energy are not unique.
Buckmaster, Colombo and Vicol~\cite{bcv2022} constructed non-unique weak
solutions that are smooth outside a set of singular times of Hausdorff
dimension less than one, and Cheskidov and Luo~\cite{cheskidov-luo2022}
proved non-uniqueness on the torus in $L^p_tL^\infty_x$ for every $p<2$,
which is sharp in view of the Ladyzhenskaya--Prodi--Serrin criterion; we
refer to the survey~\cite{buckmaster-vicol2021}. These solutions do not satisfy the energy inequality, so they are not
Leray--Hopf solutions, and convex integration leaves the question open for
Leray--Hopf solutions.

The second method is instability. In self-similar variables, an unstable
eigenvalue of the linearization around a self-similar solution becomes,
after the change back to physical variables, a second solution with the
same datum. Vishik~\cite{vishik2018a,vishik2018b} carried this out for the
two-dimensional Euler equations with vorticity in $L^p$, with a force in
$L^1_t(L^1\cap L^p)_x$; a complete account is given by Albritton, Bru\'e,
Colombo, De~Lellis, Giri, Janisch and Kwon~\cite{abcdgjk2024}. Albritton,
Bru\'e and Colombo~\cite{abc2022} adapted the method to the
three-dimensional Navier--Stokes equations and proved non-uniqueness of
Leray--Hopf solutions with a force; their unstable profile is a smooth
compactly supported vortex ring whose cross-section is a modification of
Vishik's unstable vortex. Castro, Faraco, Mengual and
Solera~\cite{cfms2026} obtained sharp non-uniqueness with forcing for the
$\alpha$-SQG equations, which include the Euler equations, using smooth
compactly supported unstable vortices. In all of these results the force is chosen so that an unstable profile
is an exact solution of the forced equation. Without force, the
background must be an exact solution, and constructing an exact solution
whose linearization is provably unstable is the main difficulty; see the
discussion in~\cite{coiculescu2026a}. For the unforced two-dimensional
Euler equations, Coiculescu~\cite{coiculescu2026b} constructed non-unique
solutions with vorticity in $L^1\cap L^p$ by perturbing a stationary
radial solution, with the spectral input verified by rigorous interval
arithmetic. Choi~\cite{choi2026} constructed two solutions with the same
initial vorticity in every finite $L^p$, using a cascade of quadrupoles
across infinitely many scales. For the Navier--Stokes
equations, Coiculescu and Palasek~\cite{coiculescu-palasek2026}
constructed data in $BMO^{-1}$ with two distinct global solutions, both
smooth at positive times, which shows that the small-data theorem of Koch
and Tataru is sharp.

For the unforced Navier--Stokes equations, Jia and
\v Sver\'ak~\cite{jia-sverak2015} proposed the following program. Take a
scale-invariant datum and a forward self-similar solution emanating from
it, which exists by~\cite{jia-sverak2014}. If the linearization around
this solution in self-similar variables has an unstable eigenvalue, then,
under a further spectral condition, there are two distinct Leray--Hopf
solutions with the same datum, and a localization of the datum makes it
of finite energy. The spectral condition concerns a non-self-adjoint
operator around a large solution that is not explicit. Guillod and
\v Sver\'ak~\cite{guillod-sverak2023} computed self-similar solutions for
a one-parameter family of axisymmetric scale-invariant data and found a
symmetry-breaking bifurcation at large data, which would give
non-uniqueness through this program.

Ionescu, Jia and Palasek~\cite{ijp2026} constructed
numerically, with pointwise residuals of order $10^{-10}$, unstable
self-similar solutions in the axisymmetric class without swirl. Binz and
Coiculescu~\cite{binz-coiculescu2026} proved that, for sufficiently
regular data, there are no nontrivial homothetic self-similar solutions in
three dimensions.

Theorem~A carries out a version of the program of Jia and \v Sver\'ak in
which the self-similar solution is replaced by a discretely self-similar
one and the unstable eigenvalue by a Floquet exponent; both are
constructed here. The proof uses no computer assistance, and the solutions
have no swirl.

\subsection{The main result}
\label{sec:intro-result}

Write $L^2_{\rm div}=\{u\in L^2(\R^3;\R^3):\operatorname{div}u=0\}$ and
$H^k_{\rm div}=H^k\cap L^2_{\rm div}$. The weak Lebesgue space
$L^{3,\infty}$ has the quasinorm
$\|f\|_{3,\infty}=\sup_{s>0}s|\{x:|f(x)|>s\}|^{1/3}$.
We include the initial condition in the weak formulation: for a viscosity
$\nu>0$, a pair $(u,p)$ solves the equation with datum $u_0$ on $[0,T)$ if,
for every real $\varphi\in C_c^\infty(\R^3\times[0,T);\R^3)$,
\begin{equation}\label{eq:initial-weak-form}
 \int_0^T\!\int_{\R^3}
 \bigl(u\cdot\partial_t\varphi+(u\otimes u):\nabla\varphi
       +\nu u\cdot\Delta\varphi+p\operatorname{div}\varphi\bigr)\,dx\,dt
 +\int_{\R^3}u_0\cdot\varphi(x,0)\,dx=0.
\end{equation}
Equation~\eqref{eq:ns} is the case $\nu=1$; a general viscosity is used
for intermediate constructions.
The main theorem of this paper is the following.

\begin{namedtheorem}{A}\label{thm:finite-energy-target}
There are
\(T>0\), a nonzero compactly supported datum \(u_0\in L^2_{\rm div}\),
and real solutions \((u^{(\alpha)},p^{(\alpha)})\), \(0\le\alpha\le1\),
of the unforced equation with viscosity one and the same datum. The datum
is smooth away from zero and satisfies
\[
 u_0\in L^{3,\infty}(\R^3)\cap\bigcap_{1\le q<3}L^q(\R^3),
 \qquad u_0\notin L^3(\R^3).
\]
Each solution is smooth at positive times and satisfies
\[
 \begin{aligned}
 u^{(\alpha)}&\in C([0,T];L^2_{\rm div})\cap L^2(0,T;H^1),\\
 u^{(\alpha)}&\in C([\delta,T];H^2_{\rm div})\quad(0<\delta<T),\qquad
 p^{(\alpha)}\in L^{5/3}(\R^3\times(0,T)).
 \end{aligned}
\]
All solutions have the strong \(L^2\) trace \(u_0\) at zero and
satisfy \eqref{eq:initial-weak-form} with \(\nu=1\). For every
\(0\le t_0\le t_1\le T\),
\[
 \frac12\|u^{(\alpha)}(t_1)\|_2^2+
 \int_{t_0}^{t_1}\|\nabla u^{(\alpha)}(t)\|_2^2\,dt
 =\frac12\|u^{(\alpha)}(t_0)\|_2^2.
\]
They also satisfy the local energy equality. Writing
$u=u^{(\alpha)}$, $p=p^{(\alpha)}$ and $e=|u|^2/2$, for every real
$\phi\in C_c^\infty(\R^3\times[0,T))$,
\begin{equation}\label{short-loc:local-energy}
 \begin{aligned}
 \int_0^T\!\int_{\R^3}|\nabla u|^2\phi\,dx\,dt
 ={}&\int_0^T\!\int_{\R^3}
       \bigl[e(\partial_t\phi+\Delta\phi)
                    +(e+p)u\cdot\nabla\phi\bigr]\,dx\,dt\\
 &+\frac12\int_{\R^3}|u_0|^2\phi(x,0)\,dx.
 \end{aligned}
\end{equation}
In particular, the solutions are suitable. Moreover,
\[
 \sup_{0\le\alpha\le1}\|u^{(\alpha)}\|_{L^\infty(0,T;L^{3,\infty})}<\infty,
 \qquad
 \sup_{0\le\alpha\le1}\|u^{(\alpha)}\|_{L^r(0,T;L^q)}<\infty
\]
for every $1\le r\le\infty$, $2\le q\le\infty$ with $2/r+3/q>1$.
There are constants $0<c_*\le C_*<\infty$ and $0<\sigma<1/4$ such that,
for every $\alpha,\alpha'\in[0,1]$ and $0<t\le T$,
\[
 c_*|\alpha-\alpha'|t^{\sigma+1/4}
 \le\|u^{(\alpha)}(t)-u^{(\alpha')}(t)\|_2
 \le C_*|\alpha-\alpha'|t^{\sigma+1/4}.
\]
In particular,
\[
 \int_0^T\|u^{(\alpha)}-u^{(\alpha')}\|_2^2\,dt
 \ge\frac{c_*^2|\alpha-\alpha'|^2T^{2\sigma+3/2}}{2\sigma+3/2}.
\]
Thus the family consists of pairwise distinct Leray--Hopf solutions.
The datum and all solutions are axisymmetric without swirl: they are
invariant under rotations about the \(x_3\)-axis, and their azimuthal
components vanish.
\end{namedtheorem}

The integrability in Theorem~A is sharp on the scale of mixed Lebesgue
spaces. Every solution of the family is bounded in $L^r(0,T;L^q)$ for
all $1\le r\le\infty$ and $2\le q\le\infty$ with $2/r+3/q>1$. By the
weak--strong uniqueness theorem~\cite{prodi1959,serrin1963,ladyzhenskaya1967},
none of them lies in $L^r(0,T;L^q)$ with $2/r+3/q=1$ and
$3<q\le\infty$, since otherwise all of them would coincide. Similarly,
the uniform bound in $L^\infty(0,T;L^{3,\infty})$ cannot be improved to
$L^\infty(0,T;L^3)$: the strong $L^2$ trace would then put $u_0$ in $L^3$.

The rate $t^{\sigma+1/4}$ has a simple origin. To leading order the
difference of two solutions is a multiple of
$t^{\sigma-1/2}b_*(x/\sqrt t,\log t)$, where $b_*$ is the periodic
profile of a Floquet mode with exponent $\sigma$ (see
Section~\ref{sec:intro-construction}); its $L^2$ norm is a multiple of
$t^{\sigma+1/4}$.

\subsection{The construction}
\label{sec:intro-construction}

We follow the program of Jia and \v Sver\'ak~\cite{jia-sverak2015}.
In the similarity variables $y=x/\sqrt t$ and $\tau=\log t$, write
$u(x,t)=t^{-1/2}U(x/\sqrt t,\log t)$. Equation~\eqref{eq:ns} becomes
\[
 \partial_\tau U=\Delta U+\tfrac12y\cdot\nabla U+\tfrac12U
 -\mathbb P\operatorname{div}(U\otimes U),
\]
where $\mathbb P$ is the Leray projection, and the initial time $t=0$
becomes $\tau=-\infty$. Self-similar solutions are stationary in $\tau$.
Discretely self-similar solutions, those with
$u(\lambda x,\lambda^2t)=\lambda^{-1}u(x,t)$ for one $\lambda>1$, are
periodic in $\tau$ with period $2\log\lambda$. Suppose that $U_*$ is a
periodic solution and that its linearization has a Floquet solution
$e^{\sigma\tau}b_*$, with $\sigma>0$ and $b_*$ periodic. Since
$e^{\sigma\tau}=t^\sigma$ vanishes at $t=0$, the fields
$U_*+\alpha e^{\sigma\tau}b_*$, $\alpha\in\R$, all carry the datum of
$U_*$, and one expects to correct each of them to an exact solution. The
difficulty is to find an exact unforced periodic solution whose
linearization provably has such a mode. Proposition~\ref{prop:unstable-background}
in Section~\ref{sec:reduction} states the background and the mode we
construct. The construction has four steps.

\paragraph{Step 1: the planar vortex train and its Gaussian instability
(Sections~\ref{sct:section}--\ref{sec:gaussian-evolution}).}
The building block is the planar vortex with Gaussian stream function
$\psi_0(x)=-2e^{-|x|^2/4}$, $x\in\R^2$. Its velocity
$\nabla^\perp\psi_0=e^{-|x|^2/4}x^\perp$ rotates, and its vorticity
$(2-|x|^2/2)e^{-|x|^2/4}$ changes sign and has zero integral: it is a
shielded vortex. The instability of shielded vortices is classical in
geophysical fluid dynamics~\cite{gent-mcwilliams1986}. We prove that the
Euler equations linearized around $\psi_0$ have exactly one unstable pair
of eigenvalues $\lambda,\bar\lambda$, with $\lambda=\gamma+i\omega$,
$\gamma>0$, and eigenfunctions in angular mode two
(Proposition~\ref{prop:short-gaussian-seed}). The count uses the index
theorem of Lin and Zeng~\cite{lin-zeng2021} together with an explicit
negative direction. Lin~\cite{lin2003} developed related criteria for
instability. With a small viscosity
$\nu>0$ the vortex diffuses: at time $s$ its vorticity is
$a^{-2}E_0(x/\sqrt a)$, $a=1+\nu s$, where $E_0$ is the vorticity of
$\psi_0$. The turnover clock $s/(1+\nu s)$ has total length
$\nu^{-1}$. In the heat variables the unstable evolution has the canonical
normalization $N_\nu=e^{\lambda/\nu}\nu^{-\kappa}$, for a constant
$\kappa\in\C$ (Proposition~\ref{prop:short-canonical}). A fixed canonical
age $\theta=[\nu(1+\nu s)]^{-1}>0$ corresponds to physical time
$s=(\nu^2\theta)^{-1}-\nu^{-1}$, of order $\nu^{-2}$. The train is the
exact solution of the planar Navier--Stokes equations with viscosity
$\nu$ on the cylinder $(\R/c\mathbb Z)\times\R$, $c=\nu^{-2}$, whose
initial stream function is the periodic sum
$\sum_{n\in\mathbb Z}\psi_0(q-nc,\eta)$. It exists globally by the
theorem of Gallay and Slijep\v cevi\'c~\cite{gallay-slijepcevic2015}.
Until the diffusing vortices overlap the train is close to the sum of
independently evolving vortices, and afterwards it is a decaying shear
flow plus a part that decays exponentially (Lemma~\ref{sct:train}).

\paragraph{Step 2: thin rings and the exact discretely self-similar
background (Sections~\ref{sec:ring-lift}--\ref{short-bg:section}).}
In cylindrical coordinates $(r,z)$ and for $R\ge2$, use the
logarithmic ring chart
\[
 q=R\log(r/R),\qquad \eta=Rz/r,\qquad s=R^2t/r^2.
\]
The stream function $r\psi(q,\eta,s)$ of the train defines an
axisymmetric field without swirl. The planar vortex centered at $q=nc$
becomes a vortex ring of radius $r_n=Re^{nc/R}$ whose cross-section has
size comparable to $r_n/R$; the rings accumulate at the origin and at
infinity. The dilation $(r,z,t)\mapsto(\Lambda r,\Lambda z,\Lambda^2t)$
with $\Lambda=e^{c/R}$ sends $q$ to $q+c$ and fixes $\eta$ and $s$, so
the periodicity of the train becomes exact discrete self-similarity:
translation becomes discrete scaling. The lifted field solves the
Navier--Stokes equations up to a curvature residual, which is small when
the rings are thin, $\epsilon=R^{-1}\to0$
(Proposition~\ref{sct:lift}). In similarity variables the residual is a
small periodic force. Since the background must be an exact unforced
solution, we remove it: a right inverse for the linearized periodic
problem and a nonlinear correction give an exact periodic solution
$U_*$. Its datum $u_{\rm dss}$ is discretely self-similar and of size
$|x|^{-1}$, and $U_*$ differs from the heat evolution of $u_{\rm dss}$
by a remainder in $H^2$ (Section~\ref{short-bg:section}). The three
small parameters are tied together: $c=\nu^{-2}$ and $\epsilon=e^{-d/\nu}$
with $d$ in a fixed compact interval.

\paragraph{Step 3: the direct three-dimensional Floquet construction
(Sections~\ref{sec:stress-spaces-short}--\ref{sec:spectral-assembly}).}
We fix the exact background and look for $0<\sigma<1/4$ and a nonzero
periodic profile $b_*$ solving~\eqref{loc:eq:mode}. The mode comes from
a loop through the rings. In similarity variables the rings move toward
the origin and age as $\tau$ increases: rings far from the origin are
young and still close to the inviscid Gaussian vortex, and rings near
the origin are old and have diffused. A strain acting on a young ring
excites its unstable angular-two mode, which grows by the factor
$N_\nu$ during the life of the ring. Once the ring is old, the stress of
the amplified perturbation generates, through the three-dimensional
pressure, a strain on the younger rings. By discrete self-similarity this
loop repeats in every period, and a Floquet mode is a loop that
reproduces its input up to the factor $e^{-\sigma\mathcal T}$, where
$\mathcal T$ is the period.

We find the exponent in the spirit of the Birman--Schwinger
principle~\cite{birman1961,schwinger1961}: instead of studying the
spectrum of the full linearized operator, we reduce the Floquet problem to
a scalar equation for the gain of one loop. We solve the forced
linearized equation along the train with zero initial data and bounds polynomial in
$\nu^{-1}$ (Section~\ref{sec:planar-response}), and rewrite the Floquet
problem as a fixed point problem for the return operator, which maps the
strain received by young rings to the strain they return. In the joint
limit $\epsilon,\nu\to0$, and on a fixed space of smooth sources, this
operator is a rank-one operator plus a remainder that tends to zero. Its
nonzero eigenvalue, the gain of one loop for an exponent $\mu$, is
\[
 2\epsilon^3\nu^{5+2\mu}\operatorname{Re}\{\mathcal C(\mu)N_\nu\}+o(1),
\]
where $\mathcal C$ is an explicit holomorphic function that does not
vanish identically. The factor $\epsilon^3$ measures the weak
three-dimensional coupling of thin rings, and $N_\nu$ is the planar
amplification. We choose $\nu_n\to0$ so that the phase of
$\mathcal C N_{\nu_n}$ is fixed and $\epsilon_n$ so that the gain is
close to one, construct the background once at these parameters, and vary
only $\mu$ in a window of width $|\log\nu_n|^{-1}$. Across the window the
gain moves from about $e$ to about $e^{-1}$. Its crossing of one gives a
Floquet exponent $\sigma_n$ of the exact three-dimensional operator, with
its full pressure and the correction of Step~2
(Section~\ref{sec:spectral-assembly}). The planar response only supplies
the limit of the return operator; the mode is constructed for the exact
operator.

\paragraph{Step 4: nonlinear correction and localization
(Section~\ref{sec:localization}).}
With the background and the mode in hand we follow the instability and
localization argument of Jia and \v Sver\'ak~\cite{jia-sverak2015},
adapted to a profile periodic in $\tau$. We replace the infinite-energy
datum $u_{\rm dss}$ by a compactly supported divergence-free truncation
that equals $u_{\rm dss}$ near the origin. The resulting common heat
correction $K$ satisfies $\|K(\tau)\|_{W^{2,\infty}}=O(e^{\tau/2})$.
For each amplitude $\alpha\in[0,1]$ we seek a profile
$U_*+K+\alpha e^{\sigma\tau}b_*+z_\alpha$ with
$\|z_\alpha(\tau)\|_{H^2}=O(e^{\beta\tau})$ as $\tau\to-\infty$,
where $\sigma<\beta<2\sigma<1/2$. We choose $\beta$ so that no Floquet
multiplier has modulus $e^{\beta\mathcal T}$ and split the periodic
evolution at this rate. The resulting right inverse accounts for every
faster unstable mode, so no maximal-growth assumption on $\sigma$ is
needed. Both forcing rates $2\sigma$ and $1/2$ exceed $\beta$, and a
contraction gives the nonlinear corrections. Their Lipschitz dependence
on $\alpha$, uniformly on one interval $[0,T]$, gives the two-sided
separation estimate of Theorem~A.

\paragraph{Organization.}
Section~\ref{sec:reduction} fixes conventions and states
Proposition~\ref{prop:unstable-background}. Sections~\ref{sct:section}--\ref{sec:localization} carry out the four steps above, and the proof of
Theorem~A is completed at the end of Section~\ref{sec:localization}.

\paragraph{Formalization.}
A Lean formalization of this work is forthcoming at
\url{https://github.com/mcoiculescu/unforced-leray-hopf-public}.

\subsection*{Declaration on the use of large language models}
\addcontentsline{toc}{subsection}{Declaration on the use of large language models}

This work was carried out with extensive use of the large language models
Claude (Anthropic) and Codex (OpenAI). The models were used to search for
the construction, to write and revise the proofs, and to check them.

While exploring the instability of an exact self-similar profile as a
mechanism, among others, for nonuniqueness, GPT-6 Astra proposed using a
planar unstable vortex as a building block to lift to an unstable
profile. We found this approach promising and pushed the model to
complete the construction of an exact discretely self-similar solution
with a growing mode.

The argument first appeared in a draft of about 200 pages, written largely
by the models, whose exposition was poor and highly redundant. We then
realized that the construction does not need swirl; working in the
axisymmetric class without swirl reduced the draft to about 150 pages.
Finally, we replaced the original proof of the existence of the Floquet
exponent by a more standard argument in the spirit of the Birman--Schwinger
principle, which reduces the Floquet problem to a scalar equation for the
gain of one loop (Section~\ref{sec:intro-construction}). This cut the
length of the proof in half, to about 70 pages.

The authors take full responsibility for the content of this paper.

\section{Conventions and the reduction to a growing mode}
\label{sec:reduction}

We use $(u\otimes v)_{ij}=u_iv_j$,
$(\operatorname{div}T)_i=\partial_jT_{ij}$,
$R_j=\partial_j(-\Delta)^{-1/2}$ and the Cartesian Leray projection
$\mathfrak P=\mathbb P$. The pressure is the global choice
$p=\sum_{i,k}R_iR_k(u_i u_k)$. A vector field is axisymmetric without
swirl if it is invariant under rotations about the $x_3$-axis and its
azimuthal component vanishes.
For intermediate constructions the viscosity is $\nu>0$, and the initial
condition is included through~\eqref{eq:initial-weak-form}.

The construction uses the similarity variables $y=x/\sqrt t$,
$\tau=\log t$. For a period $\mathcal T>0$, set
\[
 \mathcal Y_q=L^q([0,\mathcal T],d\tau/\mathcal T;L^q(\mathbb R^3)),
 \qquad \mathcal T_0=\partial_\tau-\tfrac12y\cdot\nabla.
\]
Periodic spaces identify their endpoints. The analytic construction
needed for localization is stated below.

\begin{proposition}\label{prop:unstable-background}
There exist \(\nu>0\), \(\mathcal T>0\), and a nonzero real divergence-free
field \(u_{\rm dss}\) on \(\mathbb R^3\), smooth away from zero, satisfying
\begin{equation}
 u_{\rm dss}(e^{\mathcal T/2}x)=e^{-\mathcal T/2}u_{\rm dss}(x),\qquad
 |D^\alpha u_{\rm dss}(x)|\le C_\alpha |x|^{-1-|\alpha|}
 \quad (x\ne0,\ |\alpha|\le4).
 \label{loc:eq:dss}
\end{equation}
Set
\[
 H_\nu[u_{\rm dss}](y,\tau)=e^{\tau/2}
       [e^{\nu e^\tau\Delta}u_{\rm dss}](e^{\tau/2}y).
\]
There is a real divergence-free \(\mathcal T\)-periodic distributional
solution \(U_*\) of
\begin{equation}
 \partial_\tau U_*=\nu\Delta U_*+\tfrac12 y\cdot\nabla U_*+\tfrac12 U_*-\mathfrak P\operatorname{div}(U_*\otimes U_*),
 \qquad U_*=H_\nu[u_{\rm dss}]+r_{\rm heat},\quad r_{\rm heat}\in C_{\rm per}H^2(\mathbb R^3).
 \label{loc:eq:background}
\end{equation}

There are \(0<\sigma<1/4\) and a nonzero real \(\mathcal T\)-periodic
divergence-free field \(b_*\) satisfying
\begin{equation}
 b_*\in\mathcal Y_{2,\rm div},\qquad \nabla b_*\in\mathcal Y_2.
 \label{loc:eq:energy-mode}
\end{equation}
The second space is tensor-valued. The periodic field \(b_*\) is the
Floquet profile: \(e^{\sigma\tau}b_*\) solves the linearized equation
and has multiplier \(e^{\sigma\mathcal T}>1\) over one period.
The equation for the profile is
\begin{equation}
 \begin{gathered}
 \sigma b_*+\mathcal T_0 b_*-\nu\Delta b_*-\tfrac12 b_*
 +\mathfrak P\operatorname{div}(U_*\otimes b_*+b_*\otimes U_*)=0\\
 \hbox{in }\mathcal D'(\R^3\times\R/\mathcal T\mathbb Z).
 \end{gathered}
 \label{loc:eq:mode}
\end{equation}
The fields \(u_{\rm dss}\), \(U_*\) and \(b_*\) are axisymmetric
without swirl.

\end{proposition}

Proposition~\ref{short-loc:reduction} in Section~\ref{sec:localization}
shows that Proposition~\ref{prop:unstable-background} implies Theorem~A.
Constants denoted by $P_\nu$ or $\mathcal P_\nu$ are fixed polynomials
in $\nu^{-1}$ and the specified polynomial cell parameters. Each
exponential loss is kept separately until the final parameter
inequalities are imposed.

\ProofModule{seed-train/train.tex}{The exact periodic planar train and
its complete velocity, pressure and heat-tail estimates are needed.}

\section{The Gaussian instability and canonical evolution}
\label{sec:gaussian-evolution}
\ProofModule{seed-train/lin-seed.tex}{The Gaussian seed, its direct and
adjoint couplings, and the complementary Euler evolution are needed.}
\ProofModule{seed-train/canonical.tex}{The canonical heat branch
and nonzero Mellin coupling are needed.}

\ProofModule{seed-train/ring-lift.tex}{The divergence-free ring lift, discrete
scaling, small residual and inlet bounds are needed.}

\section{The common free equation and local comparison}
\ProofModule{free-foundation.tex}{The common Cartesian body and stress
inverses, maximal domains, phase traces and adjoint graph are needed.}
\ProofModule{near/free-comparison.tex}{The global-input compact-output
Hilbert body and tensor comparisons are needed.}
\ProofModule{background/free-comparison.tex}{The common-measure body
$L^{11/4}$ comparison is needed.}
\ProofModule{background/short-background.tex}{The ordinary mixed solve,
exact unforced background and common heat remainder are needed.}
\ProofModule{stress-spaces.tex}{The split thin/repair stress spaces,
their exact coefficient multipliers and complete age-moment bounds
are needed.}

\section{Planar amplification, the full age tail, and the causal inverse}
\label{sec:planar-response}
\ProofModule{near/full-clock.tex}{The full actual-cylinder propagator
and its all-clock polynomial gain bounds are needed.}
\ProofModule{tails/canonical-transfer.tex}{The actual canonical column,
all-age kinetic and moment transfer, Mellin limit and stress defects
are needed.}
\ProofModule{near/flat-causal.tex}{The direct zero-entrance flat
weak-stress solver with polynomial gain bounds is needed.}
\ProofModule{separated-kernels.tex}{The split stress spaces,
multipliers, complete separated Cartesian kernels and moment bounds
are needed.}
\ProofModule{near/equation-comparison.tex}{The exact direct near
inverse and its physical forced-column comparison are needed.}

\ProofModule{ward-identity.tex}{The complete physical translation
identity, pressure residual, shifted-resolvent cutoff formula and
nilpotent factor are needed.}
\ProofModule{tails/young-response.tex}{The tested response from young
forcing to the complete old stress is needed.}
\ProofModule{return/flat-return.tex}{The signed reference return and
its nonzero canonical coefficient are needed.}
\ProofModule{return/source-class.tex}{Uniform return estimates on
one fixed smooth-source space are needed.}
\ProofModule{return/physical-affine.tex}{The exact physical affine
force, including its hoop component and cutoff terms, is needed.}
\ProofModule{repair/proposition.tex}{The repair, pressure, transitions,
other age rectangles and repeated-error bounds are needed.}
\ProofModule{mixed-return-lemma.tex}{The real two-norm return lemma
is needed.}
\ProofModule{spectral-assembly.tex}{The exact return reduction,
parameter selection, continuity, real root and energy-mode
reconstruction are needed.}
\ProofModule{localization.tex}{The nonlinear construction and
finite-energy localization proving Theorem~A are needed.}

\begin{thebibliography}{99}
\small

\bibitem{abc2022}
Dallas Albritton, Elia Bru\'{e}, and Maria Colombo, Non-uniqueness of Leray solutions of the forced
Navier--Stokes equations, \emph{Ann. of Math. (2)} \textbf{196} (2022), no.~1, 415--455.
\href{https://doi.org/10.4007/annals.2022.196.1.3}{doi:10.4007/annals.2022.196.1.3}

\bibitem{abcdgjk2024}
Dallas Albritton, Elia Bru\'{e}, Maria Colombo, Camillo De Lellis, Vikram Giri, Maximilian Janisch, and
Hyunju Kwon, \emph{Instability and Non-uniqueness for the 2D Euler Equations, after M.~Vishik},
Annals of Mathematics Studies, vol.~219, Princeton University Press, Princeton, 2024.
\href{https://doi.org/10.1515/9780691257846}{doi:10.1515/9780691257846}

\bibitem{bss2018}
T. Barker, G. Seregin, and V. \v{S}ver\'{a}k, On stability of weak Navier--Stokes solutions with large
$L^{3,\infty}$ initial data, \emph{Comm. Partial Differential Equations} \textbf{43} (2018), no.~4,
628--651.
\href{https://doi.org/10.1080/03605302.2018.1449219}{doi:10.1080/03605302.2018.1449219}

\bibitem{binz-coiculescu2026}
Tim Binz and Matei P. Coiculescu, Homothetic self-similar solutions to the incompressible Navier--Stokes
equations, preprint, 2026.
\href{https://arxiv.org/abs/2607.12159}{arXiv:2607.12159}

\bibitem{birman1961}
M.~Sh. Birman, On the spectrum of singular boundary-value problems, \emph{Mat. Sb. (N.S.)}
\textbf{55(97)} (1961), no.~2, 125--174 (Russian); English transl.: The spectrum of singular boundary
problems, \emph{Amer. Math. Soc. Transl. (2)} \textbf{53} (1966), 23--80.
\href{https://doi.org/10.1090/trans2/053/03}{doi:10.1090/trans2/053/03}

\bibitem{bradshaw-tsai2017}
Zachary Bradshaw and Tai-Peng Tsai, Forward discretely self-similar solutions of the Navier--Stokes
equations II, \emph{Ann. Henri Poincar\'{e}} \textbf{18} (2017), no.~3, 1095--1119.
\href{https://doi.org/10.1007/s00023-016-0519-0}{doi:10.1007/s00023-016-0519-0}

\bibitem{bcv2022}
Tristan Buckmaster, Maria Colombo, and Vlad Vicol, Wild solutions of the Navier--Stokes equations whose
singular sets in time have Hausdorff dimension strictly less than 1, \emph{J. Eur. Math. Soc. (JEMS)}
\textbf{24} (2022), no.~9, 3333--3378.
\href{https://doi.org/10.4171/JEMS/1162}{doi:10.4171/JEMS/1162}

\bibitem{buckmaster-vicol2019}
Tristan Buckmaster and Vlad Vicol, Nonuniqueness of weak solutions to the Navier--Stokes equation,
\emph{Ann. of Math. (2)} \textbf{189} (2019), no.~1, 101--144.
\href{https://doi.org/10.4007/annals.2019.189.1.3}{doi:10.4007/annals.2019.189.1.3}

\bibitem{buckmaster-vicol2021}
Tristan Buckmaster and Vlad Vicol, Convex integration constructions in hydrodynamics,
\emph{Bull. Amer. Math. Soc. (N.S.)} \textbf{58} (2021), no.~1, 1--44.
\href{https://doi.org/10.1090/bull/1713}{doi:10.1090/bull/1713}

\bibitem{ckn1982}
L. Caffarelli, R. Kohn, and L. Nirenberg, Partial regularity of suitable weak solutions of the
Navier--Stokes equations, \emph{Comm. Pure Appl. Math.} \textbf{35} (1982), no.~6, 771--831.
\href{https://doi.org/10.1002/cpa.3160350604}{doi:10.1002/cpa.3160350604}

\bibitem{cfms2026}
\'{A}ngel Castro, Daniel Faraco, Francisco Mengual, and Marcos Solera, Unstable vortices, sharp
non-uniqueness with forcing, and global smooth solutions for the SQG equation,
\emph{Arch. Ration. Mech. Anal.} \textbf{250} (2026), no.~6, Paper No.~97.
\href{https://doi.org/10.1007/s00205-026-02252-x}{doi:10.1007/s00205-026-02252-x}

\bibitem{chae-wolf2018}
Dongho Chae and J\"{o}rg Wolf, Existence of discretely self-similar solutions to the Navier--Stokes
equations for initial value in $L^2_{\mathrm{loc}}(\mathbb{R}^3)$, \emph{Ann. Inst. H. Poincar\'{e} C
Anal. Non Lin\'{e}aire} \textbf{35} (2018), no.~4, 1019--1039.
\href{https://doi.org/10.1016/j.anihpc.2017.10.001}{doi:10.1016/j.anihpc.2017.10.001}

\bibitem{cheskidov-luo2022}
Alexey Cheskidov and Xiaoyutao Luo, Sharp nonuniqueness for the Navier--Stokes equations,
\emph{Invent. Math.} \textbf{229} (2022), no.~3, 987--1054.
\href{https://doi.org/10.1007/s00222-022-01116-x}{doi:10.1007/s00222-022-01116-x}

\bibitem{choi2026}
Hyungjun Choi, Nonuniqueness of solutions to 2D incompressible Euler
equations with unbounded vorticity, preprint, 2026.
\href{https://arxiv.org/abs/2610.37653}{arXiv:2610.37653}

\bibitem{coiculescu2026a}
Matei P. Coiculescu, Ill-posedness of the Euler equations linearized around
homogeneous steady states, preprint, 2026.
\href{https://arxiv.org/abs/2609.06313}{arXiv:2609.06313}

\bibitem{coiculescu2026b}
Matei P. Coiculescu, Non-uniqueness of weak solutions to the unforced
two-dimensional Euler equations with vorticity in $L^1\cap L^p$, I: A
compactness criterion and the linear analysis, preprint, 2026.
\href{https://arxiv.org/abs/2610.30655}{arXiv:2610.30655}

\bibitem{coiculescu-palasek2026}
Matei P. Coiculescu and Stan Palasek, Non-uniqueness of smooth solutions of the Navier--Stokes equations
from critical data, \emph{Invent. Math.} \textbf{244} (2026), no.~1, 165--219.
\href{https://doi.org/10.1007/s00222-025-01396-z}{doi:10.1007/s00222-025-01396-z}

\bibitem{engel-nagel2000}
Klaus-Jochen Engel and Rainer Nagel,
\emph{One-Parameter Semigroups for Linear Evolution Equations},
Graduate Texts in Mathematics, vol.~194, Springer-Verlag,
New York, 2000.
\href{https://doi.org/10.1007/b97696}{doi:10.1007/b97696}.
\href{https://www.math.uni-tuebingen.de/de/forschung/agfa/members/engel-nagel_one-parameter-semigroups.pdf}{Authors' online edition}.

\bibitem{ess2003}
Luis Escauriaza, Grigorii A. Seregin, and Vladim\'{\i}r \v{S}ver\'{a}k, $L_{3,\infty}$-solutions of the
Navier--Stokes equations and backward uniqueness, \emph{Russian Math. Surveys} \textbf{58} (2003),
no.~2, 211--250.
\href{https://doi.org/10.1070/RM2003v058n02ABEH000609}{doi:10.1070/RM2003v058n02ABEH000609}
(Translation of \emph{Uspekhi Mat. Nauk} \textbf{58} (2003), no.~2, 3--44,
\href{https://doi.org/10.4213/rm609}{doi:10.4213/rm609}.)

\bibitem{feng-sverak2015}
Hao Feng and Vladim\'{\i}r \v{S}ver\'{a}k, On the Cauchy problem for axi-symmetric vortex rings,
\emph{Arch. Ration. Mech. Anal.} \textbf{215} (2015), no.~1, 89--123.
\href{https://doi.org/10.1007/s00205-014-0775-4}{doi:10.1007/s00205-014-0775-4}

\bibitem{gallay-slijepcevic2015}
Thierry Gallay and Sini\v{s}a Slijep\v{c}evi\'{c},
Uniform boundedness and long-time asymptotics for the two-dimensional
Navier--Stokes equations in an infinite cylinder,
\emph{Journal of Mathematical Fluid Mechanics} \textbf{17} (2015),
no.~1, 23--46.
\href{https://doi.org/10.1007/s00021-014-0188-z}{doi:10.1007/s00021-014-0188-z}.

\bibitem{gallay-sverak2019}
Thierry Gallay and Vladim\'{\i}r \v{S}ver\'{a}k, Uniqueness of axisymmetric viscous flows originating
from circular vortex filaments, \emph{Ann. Sci. \'{E}c. Norm. Sup\'{e}r.} \textbf{52} (2019), no.~4,
1025--1071.
\href{https://doi.org/10.24033/asens.2402}{doi:10.24033/asens.2402}

\bibitem{gent-mcwilliams1986}
Peter R. Gent and James C. McWilliams, The instability of barotropic circular vortices,
\emph{Geophys. Astrophys. Fluid Dynam.} \textbf{35} (1986), no.~1--4, 209--233.
\href{https://doi.org/10.1080/03091928608245893}{doi:10.1080/03091928608245893}

\bibitem{guillod-sverak2023}
Julien Guillod and Vladim\'{\i}r \v{S}ver\'{a}k, Numerical investigations of non-uniqueness for the
Navier--Stokes initial value problem in borderline spaces, \emph{J. Math. Fluid Mech.} \textbf{25}
(2023), no.~3, Paper No.~46.
\href{https://doi.org/10.1007/s00021-023-00789-5}{doi:10.1007/s00021-023-00789-5}

\bibitem{hopf1951}
Eberhard Hopf, \"{U}ber die Anfangswertaufgabe f\"{u}r die hydrodynamischen Grundgleichungen,
\emph{Math. Nachr.} \textbf{4} (1951), 213--231.
\href{https://doi.org/10.1002/mana.3210040121}{doi:10.1002/mana.3210040121}

\bibitem{ijp2026}
Alexandru D. Ionescu, Hao Jia, and Stan Palasek, On the non-uniqueness of solutions of the axi-symmetric
swirl-free Navier--Stokes equations, I, preprint, 2026.
\href{https://arxiv.org/abs/2606.07501}{arXiv:2606.07501}

\bibitem{jia-sverak2014}
Hao Jia and Vladim\'{\i}r \v{S}ver\'{a}k, Local-in-space estimates near initial time for weak solutions of
the Navier--Stokes equations and forward self-similar solutions, \emph{Invent. Math.} \textbf{196}
(2014), no.~1, 233--265.
\href{https://doi.org/10.1007/s00222-013-0468-x}{doi:10.1007/s00222-013-0468-x}

\bibitem{jia-sverak2015}
Hao Jia and Vladim\'{i}r \v{S}ver\'{a}k,
Are the incompressible 3d Navier--Stokes equations locally ill-posed
in the natural energy space?,
\emph{Journal of Functional Analysis} \textbf{268} (2015),
no.~12, 3734--3766.
\href{https://doi.org/10.1016/j.jfa.2015.04.006}{doi:10.1016/j.jfa.2015.04.006}.
The instability and localization mechanism is described in
\href{https://arxiv.org/abs/1306.2136v1}{arXiv:1306.2136v1}, Sections~4--5;
the periodic version used here is proved in Section~\ref{sec:localization}.

\bibitem{kato1984}
Tosio Kato, Strong $L^p$-solutions of the Navier--Stokes equation in $\mathbb{R}^m$, with applications
to weak solutions, \emph{Math. Z.} \textbf{187} (1984), no.~4, 471--480.
\href{https://doi.org/10.1007/BF01174182}{doi:10.1007/BF01174182}

\bibitem{koch-tataru2001}
Herbert Koch and Daniel Tataru, Well-posedness for the Navier--Stokes equations,
\emph{Adv. Math.} \textbf{157} (2001), no.~1, 22--35.
\href{https://doi.org/10.1006/aima.2000.1937}{doi:10.1006/aima.2000.1937}

\bibitem{kozono-sohr1996}
Hideo Kozono and Hermann Sohr, Remark on uniqueness of weak solutions to the Navier--Stokes equations,
\emph{Analysis} \textbf{16} (1996), no.~3, 255--272.
\href{https://doi.org/10.1524/anly.1996.16.3.255}{doi:10.1524/anly.1996.16.3.255}

\bibitem{ladyzhenskaya1967}
O.~A. Ladyzhenskaya, Uniqueness and smoothness of generalized solutions of the Navier--Stokes equation,
\emph{Zap. Nauchn. Sem. Leningrad. Otdel. Mat. Inst. Steklov. (LOMI)} \textbf{5} (1967), 169--185 (Russian);
English transl.: \emph{Sem. Math. V. A. Steklov Math. Inst. Leningrad} \textbf{5} (1969), 60--66.

\bibitem{ladyzhenskaya1968}
O.~A. Ladyzhenskaya, Unique solvability in the large of three-dimensional Cauchy problem for the
Navier--Stokes equations in the presence of axial symmetry,
\emph{Zap. Nauchn. Sem. Leningrad. Otdel. Mat. Inst. Steklov. (LOMI)} \textbf{7} (1968), 155--177 (Russian);
English transl.: \emph{Sem. Math. V. A. Steklov Math. Inst. Leningrad} \textbf{7}, 70--79.

\bibitem{lemarie2016}
Pierre Gilles Lemari\'{e}-Rieusset, \emph{The Navier--Stokes Problem in the 21st Century},
CRC Press, Boca Raton, FL, 2016. xxii+718 pp. ISBN 978-1-4665-6621-7.
\href{https://doi.org/10.1201/b19556}{doi:10.1201/b19556}

\bibitem{lmnp1999}
S. Leonardi, Josef M\'{a}lek, Jind\v{r}ich Ne\v{c}as, and Milan Pokorn\'{y}, On axially symmetric
flows in $\mathbb{R}^3$, \emph{Z. Anal. Anwend.} \textbf{18} (1999), no.~3, 639--649.
\href{https://doi.org/10.4171/ZAA/903}{doi:10.4171/ZAA/903}

\bibitem{leray1934}
Jean Leray, Sur le mouvement d'un liquide visqueux emplissant l'espace,
\emph{Acta Math.} \textbf{63} (1934), 193--248.
\href{https://doi.org/10.1007/BF02547354}{doi:10.1007/BF02547354}

\bibitem{lin2003}
Zhiwu Lin, Instability of some ideal plane flows, \emph{SIAM J. Math. Anal.} \textbf{35} (2003), no.~2,
318--356.
\href{https://doi.org/10.1137/S0036141002406266}{doi:10.1137/S0036141002406266}

\bibitem{lin-zeng2021}
Zhiwu Lin and Chongchun Zeng,
\emph{Instability, index theorem, and exponential trichotomy for
Linear Hamiltonian PDEs},
\emph{Memoirs of the American Mathematical Society} \textbf{275} (2022),
no.~1347. \href{https://doi.org/10.1090/memo/1347}{doi:10.1090/memo/1347}.

\bibitem{prodi1959}
Giovanni Prodi, Un teorema di unicit\`{a} per le equazioni di Navier--Stokes,
\emph{Ann. Mat. Pura Appl. (4)} \textbf{48} (1959), 173--182.
\href{https://doi.org/10.1007/BF02410664}{doi:10.1007/BF02410664}

\bibitem{scheffer1976}
Vladimir Scheffer, Partial regularity of solutions to the Navier--Stokes equations,
\emph{Pacific J. Math.} \textbf{66} (1976), no.~2, 535--552.
\href{https://doi.org/10.2140/pjm.1976.66.535}{doi:10.2140/pjm.1976.66.535}

\bibitem{schwinger1961}
Julian Schwinger, On the bound states of a given potential, \emph{Proc. Natl. Acad. Sci. USA}
\textbf{47} (1961), no.~1, 122--129.
\href{https://doi.org/10.1073/pnas.47.1.122}{doi:10.1073/pnas.47.1.122}

\bibitem{seregin-sverak2017}
G. Seregin and V. \v{S}ver\'{a}k, On global weak solutions to the Cauchy problem for the Navier--Stokes
equations with large $L_3$-initial data, \emph{Nonlinear Anal.} \textbf{154} (2017), 269--296.
\href{https://doi.org/10.1016/j.na.2016.01.018}{doi:10.1016/j.na.2016.01.018}

\bibitem{serrin1963}
James Serrin, The initial value problem for the Navier--Stokes equations,
in: \emph{Nonlinear Problems} (Proc. Sympos., Madison, Wis., 1962), R.~E.~Langer (ed.),
University of Wisconsin Press, Madison, 1963, pp.~69--98.

\bibitem{tsai2014}
Tai-Peng Tsai, Forward discretely self-similar solutions of the Navier--Stokes equations,
\emph{Comm. Math. Phys.} \textbf{328} (2014), no.~1, 29--44.
\href{https://doi.org/10.1007/s00220-014-1984-2}{doi:10.1007/s00220-014-1984-2}

\bibitem{ukhovskii-yudovich1968}
M.~R. Ukhovskii and V.~I. Yudovich, Axially symmetric flows of ideal and viscous fluids filling the whole space,
\emph{J. Appl. Math. Mech.} \textbf{32} (1968), no.~1, 52--62.
\href{https://doi.org/10.1016/0021-8928(68)90147-0}{doi:10.1016/0021-8928(68)90147-0}
(Translation of \emph{Prikl. Mat. Mekh.} \textbf{32} (1968), 59--69.)

\bibitem{vishik2018a}
Misha Vishik, Instability and non-uniqueness in the Cauchy problem for the Euler equations of an ideal
incompressible fluid. Part I, preprint, 2018.
\href{https://arxiv.org/abs/1805.09426}{arXiv:1805.09426}

\bibitem{vishik2018b}
Misha Vishik, Instability and non-uniqueness in the Cauchy problem for the Euler equations of an ideal
incompressible fluid. Part II, preprint, 2018.
\href{https://arxiv.org/abs/1805.09440}{arXiv:1805.09440}

\end{thebibliography}


\end{document}
